% Required packages in the main document:
% \usepackage{array}
% \usepackage{booktabs}
% \usepackage{tabularx}
% Keep the existing citation configuration supporting \citep.

\begin{table*}[!t]
    \centering
    \caption{Representative benchmarks for interaction rhythm in spoken dialogue.}
    \label{tab:full-duplex-benchmarks}

    \begingroup
    \footnotesize
    \setlength{\tabcolsep}{3pt}
    \renewcommand{\arraystretch}{1.08}

    \begin{tabularx}{\textwidth}{@{}
        >{\raggedright\arraybackslash}p{0.17\textwidth}
        >{\raggedright\arraybackslash}p{0.20\textwidth}
        >{\raggedright\arraybackslash}p{0.13\textwidth}
        >{\raggedright\arraybackslash}p{0.10\textwidth}
        >{\raggedright\arraybackslash}X
        >{\raggedright\arraybackslash}p{0.08\textwidth}
    @{}}
        \toprule
        \textbf{Benchmark}
        & \textbf{Task}
        & \textbf{Scale}
        & \textbf{Modality}
        & \textbf{Evaluation Focus / Metrics}
        & \textbf{Source} \\
        \midrule

        \multicolumn{6}{@{}l}{
            \textit{System-Level Interaction Benchmarks}
        } \\
        \midrule

        Talking Turns
        & Turn-taking; interruptions; agent backchannels
        & 11 participants
        & Speech, text
        & Timing-based human judgments; pauses, gaps,
          and overlaps in human--system conversations.
        & \citep{DBLP:conf/iclr/AroraLCP025} \\
        \addlinespace[2pt]

        FDB-v1\textsuperscript{*}
        & Turn-taking; pauses; interruptions; agent backchannels
        & 727 samples
        & Speech, text
        & Task-dependent take-over rates; response latency;
          backchannel behavior and response quality.
        & \citep{DBLP:conf/asru/LinLLWALL25} \\
        \addlinespace[2pt]

        FD-Bench\textsuperscript{*}
        & Interruption handling; response quality
        & 293 dialogues
        & Speech, text
        & Interruption and response success; latency;
          response quality under varying noise conditions.
        & \citep{DBLP:conf/interspeech/PengCNMN0C25} \\
        \addlinespace[2pt]

        FDB-v2\textsuperscript{*}
        & Multi-turn coordination; corrections;
          task-state continuity
        & 4 task families
        & Speech, text
        & Adaptive examiner; turn-taking fluency;
          instruction following and task outcomes
          under Fast/Slow pacing.
        & \citep{DBLP:conf/acl/LinKSCAWL26} \\
        \addlinespace[2pt]

        MTR-DuplexBench
        (conversational features)
        & Multi-round turn-taking; interruptions;
          backchannels; state maintenance
        & 200 dialogues;
          10 rounds each
        & Speech, text
        & Feature-specific success rates; latency;
          backchannel frequency across multiple rounds.
        & \citep{DBLP:conf/acl/HeCXLZBMK26} \\
        \addlinespace[2pt]

        HumDial-FDBench\textsuperscript{*}
        & Interruption handling; user backchannels;
          rejection of non-addressed speech
        & 5,000 test cases
        & Speech, text;
          Chinese, English
        & Interruption and rejection accuracy;
          stopping and response latency in
          human-recorded interaction scenarios.
        & \citep{DBLP:journals/corr/abs-2601-05564} \\
        \addlinespace[2pt]

        Duplex-MPE\textsuperscript{*}
        & Multi-party participation; addressee-sensitive
          responding; silence and stopping
        & 2,000 pairs;
          4,000 streams
        & Speech, text
        & Response initiation; answer accuracy;
          silence preservation; stopping after
          request resolution.
        & \citep{DBLP:journals/corr/abs-2609-31948} \\
        \addlinespace[2pt]

        DuplexAct-Bench\textsuperscript{*}
        & Proactive initiation; silence;
          interruption handling; agent backchannels
        & 1,290 trials;
          30 scenarios
        & Speech, text;
          Chinese, English
        & Behavior-specific timing and content;
          joint behavioral correctness;
          backchannel alignment.
        & \citep{DBLP:journals/corr/abs-2609-39446} \\

        \midrule
        \multicolumn{6}{@{}l}{
            \textit{Targeted Diagnostic Benchmarks}
        } \\
        \midrule

        TurnBench
        & Turn-end and interruption detection
        & $\sim$30 h;
          154 dialogues
        & Speech, text
        & Event recall; false-positive rate;
          causal detection latency.
        & \citep{DBLP:journals/corr/abs-2608-25218} \\
        \addlinespace[2pt]

        SID-Bench\textsuperscript{*}
        & Interruption intent; rejection of
          non-interruptive feedback
        & 3,700 instances
        & Speech, text
        & False-interruption rate;
          semantic-breakpoint latency;
          average penalty time.
        & \citep{DBLP:journals/corr/abs-2603-24144} \\
        \addlinespace[2pt]

        ECHO\textsuperscript{*}
        & Context-sensitive yielding;
          backchannel and off-talk handling
        & 549 instances;
          300 pairs
        & Speech, text
        & Context-conditioned Yield/Keep decisions;
          paired action correctness and consistency.
        & \citep{DBLP:journals/corr/abs-2609-17360} \\
        \addlinespace[2pt]

        TPI-Bench
        (TPI-VA)\textsuperscript{*}
        & Third-party interruption awareness;
          response-strategy selection
        & 2,000 TPI-Test;
          2,000 Janus-Test;
          100 real-world cases
        & Speech input;
          text output
        & Response-strategy adherence;
          acoustic speaker discrimination;
          response helpfulness.
        & \citep{DBLP:conf/acl/LeeSLKY26} \\

        \bottomrule
    \end{tabularx}

    \par\vspace{3pt}
    \begin{minipage}{\textwidth}
        \footnotesize
        \raggedright
        \textit{Notes.}
        Scales refer to the listed evaluation versions or subsets.
        FDB denotes Full-Duplex-Bench and is distinct from FD-Bench.
        HumDial refers to its full-duplex track.
        \textsuperscript{*} marks links to content--rhythm coordination
        (Coordination / X), without implying a shared joint metric.
        Diagnostic results do not by themselves establish complete
        system-level interaction performance.
    \end{minipage}

    \endgroup
\end{table*}